\section{Introduction}
\label{sec:introduction}

Urban route planning is a constrained combinatorial optimization problem. A user seeking ``a one-day seafood food tour in Zhuhai'' expects a route that balances restaurant diversity, geographic proximity, time feasibility, and experiential richness---requirements that are difficult to formalize in a traditional constraint specification.

Classical approaches formulate this as a Time-Side Window Traveling Problem (TSPTW)~\cite{solomon1987algorithms}, where a set of points of interest (POIs) must be visited within their operating hours while minimizing travel time. While exact solvers produce feasible routes, they require explicit constraint definitions and cannot interpret natural language preferences such as ``romantic sunset dinner'' or ``energetic speed-run touring.''

Large Language Models (LLMs) excel at understanding user intent~\cite{wei2022chain}, but suffer from three critical limitations in spatial planning: (1)~\textit{hallucination}---recommending non-existent POIs~\cite{ji2023survey}; (2)~\textit{spatial blindness}---inability to reason about geographic distances from textual descriptions alone~\cite{gong2024spatial}; and (3)~\textit{optimism bias}---invariably attempting to satisfy impossible requests rather than declining gracefully. Recent work on LLM-based travel planning~\cite{chen2024travelagent,han2024travelplanner} has demonstrated both the potential and the limitations of single-LLM approaches: even state-of-the-art models achieve only 4.4\% success rate on multi-constraint planning tasks~\cite{han2024travelplanner}.

We observe that human travel planning naturally decomposes into specialized concerns: a food enthusiast evaluates restaurants differently from a nature lover assessing hiking trails, and both differ from a logistics expert calculating inter-POI travel times. This suggests an \textit{expert-ensemble} architecture inspired by the Mixture-of-Experts (MoE) principle~\cite{shazeer2017outrageously,jiang2024mixtral}---where domain-specific modules contribute proposals that are reviewed, refined, and synthesized into a coherent route. Unlike classical MoE systems that use learned gating networks~\cite{fedus2022switch}, our system employs rule-based scene classification for expert dispatch, providing interpretable and debuggable expert selection without requiring large-scale training data.

\subsection{Contributions}

\begin{enumerate}[leftmargin=*]
    \item \textbf{Expert-Ensemble Multi-Agent Architecture}: We design an 11-expert system with scene-aware dispatch, where experts in scenic, culinary, cultural, natural, transportation, accommodation, and local-guide domains contribute proposals in parallel waves, orchestrated by a LangGraph state machine~\cite{langgraph2024}. Unlike general-purpose multi-agent frameworks~\cite{wu2023autogen,crewai2024}, our architecture is optimized for spatial planning with dependency-ordered wave execution.

    \item \textbf{Review-Rework Feedback Loop}: A quality assurance mechanism evaluates proposals against six dimensions (intent match, POI quality, geographic continuity, scene diversity, time feasibility, category match) using both algorithmic checks and LLM judgment, and triggers targeted rework when scores fall below threshold.

    \item \textbf{Prompt Engineering for Refusal}: We demonstrate that teaching LLMs to reject infeasible requests via explicit prompt instructions---rather than architectural components---is sufficient to handle impossible scenarios, supporting the ``prompt before architecture'' design principle.

    \item \textbf{Algorithmic Time Smoothing}: We introduce \texttt{\_smooth\_times()}, a post-processing layer that detects and corrects LLM-generated temporal anomalies (gaps, inversions, overflows), embodying the ``LLM proposes, algorithm validates'' pattern.

    \item \textbf{Systematic Optimization Case Study}: We document a debugging journey from 65.2 to 81.2 points across 100 scenarios and seven rounds of targeted fixes, yielding three transferable design principles for LLM-based planning systems.
\end{enumerate}
